\documentclass{article}
\usepackage{kotex}
\usepackage{setspace}  
\setstretch{1.15}
\usepackage{xcolor}
\usepackage{graphicx}
\usepackage[linguistics]{forest}
\usepackage[a4paper, 
            left = 2.5cm,
            right = 2.5cm,
            top = 2cm,
            bottom = 2cm]{geometry}
\usepackage{fancyhdr}       
\usepackage{gb4e}
\noautomath
\glossglue=0.7em plus 0.2em minus 0.1em
\newcommand{\glabbr}[1]{\textsc{\MakeLowercase{#1}}}
\pagestyle{fancy}
\fancyhf{}                             
\fancyfoot[L]{\nouppercase{Syntax(002) Midterm}} 
\fancyfoot[R]{\thepage}                
\renewcommand{\headrulewidth}{0pt}    
\renewcommand{\footrulewidth}{0.4pt}   
\setcounter{secnumdepth}{0}

% PS-rule trees and X-bar trees
\forestset{
  terminal/.style={
    font=\itshape,
    no edge,
    l=3ex,
    for parent={l sep=0pt}
  },
  lex/.style={append={[#1, terminal]}},
  simple terminals/.style={
    delay={for tree={split option={content}{ }{content,lex}}}
  },
  ps tree/.style={
    simple terminals,
    for tree={
      align=center,
      parent anchor=south,
      child anchor=north,
      text=blue,
      edge={semithick,blue},
      inner sep=1pt,
      l sep=10mm,
      s sep=6mm
    }
  },
  xbar tree/.style={
    simple terminals,
    for tree={
      align=center,
      parent anchor=south,
      child anchor=north,
      text=blue,
      edge={semithick,blue},
      inner sep=1pt,
      l sep=8mm,
      s sep=4.5mm
    }
  }
}
\newcommand{\onebar}[1]{#1\ensuremath{^{\prime}}}
\newcommand{\twobar}[1]{#1\ensuremath{^{\prime\prime}}}

\begin{document}
\large

\title{Syntax(002) Midterm}
\author{강초롱}
\date{2025 Fall}
\maketitle
\thispagestyle{fancy}
\bigskip

\section{1. Consider the following data taken from Turkish \normalsize(3 points)}
\medskip 
\begin{exe}
  \ex
  \gll Siz \textbf{gid}-iyor.sun.uz.\\
       you.\glabbr{PL} go-\glabbr{PRES}.2.\glabbr{PL}\\
  \glt `You are going.'

  \ex
  \gll Ali \"{u}\c{c} ye\c{s}il top \textbf{al}-d\i.\\
       Ali three green ball buy-\glabbr{PST}.3\glabbr{SG}\\
  \glt `Ali bought three green balls.'

  \ex
  \gll Senin s\i{}nav\i{} \textbf{ge\c{c}}-meni
       \textbf{isti}-yorum.\\
       you test pass-\glabbr{INF} want-\glabbr{PRES}.1\glabbr{SG}\\
  \glt `I want you should pass the test.'\footnote{the matrix verb agrees with the silent matrix subject.}
\end{exe}
\bigskip
\noindent \textbf{a.} In the data, PRES means ``present tense'', PST means ``past tense''. 
What part of speech are \textbf{`gid'}, \textbf{`al'}, \textbf{`ge\c{c}'}, and \textbf{`isti'}? 
What tests did you use to figure out what speech they were? \par
\bigskip
\noindent {\color{blue} verb. \\ because they all have suffixes like tense suffix and infinitive suffix, which verbs can have.}
\vspace{1cm}

\noindent \textbf{b.} Provide the PSRs needed to generate the trees \\(assume that every clause is a CP, and \textit{\"{u}\c{c}} ``three'' is an adjective). \par
\bigskip

\noindent {\color{blue} CP → (C) TP
\\ TP → (NP) (T) VP
\\ VP → {NP/CP} V
\\ NP → (AdjP+) N
\\ AdjP → Adj} 

\newpage
\section{2. Consider the following sentence \normalsize}
\begin{center}
     \textbf{The teacher saw that her students fought constantly \\ after school in the parking lot.} \par
\end{center}
\bigskip
\noindent \textbf{a.} For the following words, indicate their part of speech \textit{AND} \par
\medskip \noindent (i) if the word is an open class part of speech, explain the (morphological/syntactic)\\ distributional criteria by which you came up with that classification.\par
\vspace{1em}
\noindent (ii) if the word is closed part of speech, indicate it with ``closed''.\par
\vspace{1em}
\noindent (iii) Mark whether each word is a lexical or functional part of speech. (4 points) \par
\vspace{2em}
\begin{description}
    \setlength{\itemsep}{2em}
    \item[\textbf{The}] {\color{blue} Determiner, closed, functional}
    \item[\textbf{Teacher}] {\color{blue} Noun, open, lexical / becuase it has noun suffix `-er'} 
    \item[\textbf{Saw}] {\color{blue} Verb, open, lexical / because it appears after the subject.}
    \item[\textbf{That}] {\color{blue} Complementizer, closed, functional} 
    \item[\textbf{Her}] {\color{blue} Determiner, closed, functional}
    \item[\textbf{Constantly}] {\color{blue} Adverb, open, lexical / because it has adverb suffix `-ly'}
    \item[\textbf{After}] {\color{blue} Preposition, closed, functional}
    \item[\textbf{Parking}] {\color{blue} Adjective, open, lexical / because it has adjective suffix `-ing' and modifies the following noun `lot' \footnote{it also can be considered as a noun.}}  
\end{description}

\newpage
\noindent \textbf{b.} The above sentence is ambiguous (assume that the ambiguity comes from the phrase ``in the parkng lot''). Draw the two PSR trees that distinguish the two possible meanings of the sentence. Provide a Korean paraphrase for each meaning. (2 points) \par
\bigskip
\noindent {\color{blue} (1) 학생들이 주차장에서 싸우는 것을 선생님이 보았다.} 
\bigskip
\begin{center}
\resizebox{\textwidth}{!}{%
\begin{forest} ps tree
[TP
  [NP
    [D The]
    [N teacher]
  ]
  [VP
    [V saw]
    [CP [C that]
        [TP 
          [NP 
            [D her]
            [N students]
          ]
          [VP 
            [V fought]
            [AdvP [Adv constantly]]
            [PP [P after] [NP [N school]]]   
            [PP [P in] [NP [D the] [AdjP [Adj parking]] [N lot]]]
            ]
]]]]
\end{forest}
}
\end{center}
\newpage
{\color{blue} (2) 학생들이 싸우는 것을 선생님이 주차장에서 보았다.} 
\bigskip
\begin{center}
\resizebox{\textwidth}{!}{%
\begin{forest} ps tree
[TP
  [NP
    [D The]
    [N teacher]
  ]
  [VP
    [V saw]
    [CP [C that]
        [TP 
          [NP 
            [D her]
            [N students]
          ]
          [VP 
            [V fought]
            [AdvP [Adv constantly]]
            [PP [P after] [NP [N school]]]]]]
    [PP [P in] [NP [D the] [AdjP [Adj parking]] [N lot]]]            
]]
\end{forest}
}
\end{center}

\newpage
\noindent \textbf{c.} Draw the two X-bar theoretic trees that distinguish the two possible meanings of the sentence. Provide a Korean paraphrase for each meaning. (2 points)
\bigskip

\noindent {\color{blue} (1) 학생들이 주차장에서 싸우는 것을 선생님이 보았다.} 
\bigskip
\begin{center}
\begin{forest} ps tree
[CP
  [\onebar{C}
    [C]
    [TP 
      [DP [\onebar{D} [D The] [NP [\onebar{N} [N teacher]]]]]
      [\onebar{T} [T] [VP [\onebar{V} [V saw] [CP
        [\onebar{C}
          [C that]
          [TP]
        ]
      ]]]]
    ]
  ]
]
\end{forest}
\end{center}
\newpage
{\color{blue} (2) 학생들이 싸우는 것을 선생님이 주차장에서 보았다.} 
\bigskip
\begin{center}
\resizebox{\textwidth}{!}{%
\begin{forest} ps tree
[TP
  [NP
    [D The]
    [N teacher]
  ]
  [VP
    [V saw]
    [CP [C that]
        [TP 
          [NP 
            [D her]
            [N students]
          ]
          [VP 
            [V fought]
            [AdvP [Adv constantly]]
            [PP [P after] [NP [N school]]]]]]
    [PP [P in] [NP [D the] [AdjP [Adj parking]] [N lot]]]            
]]
\end{forest}
}
\end{center}

\newpage
\noindent \textbf{d.} Answer the following questions (6 points) \par
\bigskip \noindent (i) Can the two theories (PSRs and X-bar theory) explain the ambiguity? If so, explain why. You must refer to the tree structures you provided above in your answer. \par
\vspace{1.5cm}
\noindent (ii) In the slightly modified sentence below, the interpretation in which \textit{the parents fought after school at the park and the teacher saw it} is available. Can both theories account for this reading? \par
\medskip \noindent \textbf{\textit{The teacher saw that her students fought constantly after school in the parking lot. She saw that their parents did so at the park as well.}}\par
\vspace{1.5cm}
\noindent (iii) Explain why (a) is ungrammatical but (b) is grammatical, referring to the rules of X-bar theory.\par
\medskip \noindent (a) \*The teacher saw \textbf{the fight} that her students fought constantly after school in the parking lot.\par
\medskip \noindent (b) The teacher saw that her students fought constantly after school in the parking lot \textbf{every day}.

\newpage
\section{3. Consider the following data from Breton, a Celtic language. \\ \normalsize{(7 points)}}
\medskip
\setcounter{exx}{0}
\begin{exe}
  \ex
  \gll ar baner c'hlas \\
       the basket blue \\
  \glt `the blue basket'

  \ex
  \gll ar vleunienn wenn \\
       the flower white \\
  \glt `the white flower'

  \ex
  \gll Kador ar bugel bihan \\
       Chair the child small \\
  \glt `The small child's chair'

  \ex
  \gll Me a-zebr pesked alies \\
       I eat fish often \\
  \glt `I often eat fish'  
\end{exe}
\bigskip
\noindent \textbf{i.} What is the setting of the specifier parameter? \par
\vspace{2em}
\noindent \textbf{ii.} What is the setting of the adjunct parameter? \par
\vspace{2em}
\noindent \textbf{iii.} What is the setting of the complement parameter? \par
\vspace{2em}
\noindent \textbf{iv.} Draw the X-bar theoretic tree for (d) and answer the following questions.
\vspace{4cm}
\begin{itemize}
     \setlength{\itemsep}{1em}
     \item[(a)] What node(s) does the VP dominate?  
     \item[(b)] What node(s) does the V head govern?
     \item[(c)] What node(s) immediately precedes AdvP?
     \item[(d)] List all the nodes that asymmetrically c-command DP \textit{fish}? 
     \item[(e)] What node(s) does the TP exhaustively dominate? 
     \item[(f)] What node(s) does the subject DP phrase govern?
\end{itemize}

\section{4. Draw the X-bar theoretic tree for the following sentence and explain why the acceptability of the following sentence challenges the binding theory we discussed in class. \normalsize{(3 points)}}
\medskip
\begin{center}
     \textbf{John's father$_i$ saw a snake near him/himself$_i$.}
\end{center}

\newpage
\section{5. Draw the X-bar theoretic tree for the following sentence, and determine whether the matrix clause and the embedded clause are finite or non-finite. You must provide evidence for your answer. \normalsize{(3 points)}}
\medskip
\begin{center} 
    \textbf{I asked him to leave.}
\end{center}

\end{document}
